\documentclass[11pt,a4paper]{article}
\usepackage[utf8]{inputenc}
\usepackage[T1]{fontenc}
\usepackage{fouriernc}
\usepackage{tikz}
\usepackage{pgfplots}
\pgfplotsset{compat=1.18}
\usepackage[french]{babel}
\frenchbsetup{StandardLists=true}
\NoAutoSpaceBeforeFDP
\usepackage{enumitem}
\usepackage{amsmath}
\usepackage{amssymb}
\usepackage[left=2cm,right=2cm,top=2cm,bottom=2cm]{geometry}
\usepackage[dvipsnames]{xcolor}
\usepackage{setspace}
\singlespacing
\usepackage{tcolorbox}
\usepackage{multirow}
\usepackage{tabularx}
\usepackage{booktabs}
\usepackage{array}
\usepackage{fancyhdr}
\setlength{\headheight}{14pt}

\pagestyle{fancy}
\renewcommand\headrulewidth{1pt}
\fancyhead[L]{Lycée Denis Diderot}
\fancyhead[R]{Classe : T~Bac Pro}
\setcounter{page}{1}
\fancyfoot[C]{Moteurs électriques --- Corrigé --- page~\thepage}
\fancyfoot[R]{\today}
\fancyfoot[L]{Alexis RUIZ}

\setlength{\parindent}{0pt}
\renewcommand{\arraystretch}{1.8}

\newcommand{\reponse}[1]{%
  \begin{tcolorbox}[colback=green!5!white, colframe=green!50!black,
                    left=2mm, right=2mm, top=1mm, bottom=1mm]
  #1
  \end{tcolorbox}}

\newcounter{nquestion}
\newcommand{\Q}[1]{%
  \stepcounter{nquestion}%
  \vspace{3mm}
  \noindent\textbf{Question \thenquestion.} \textbf{#1}
  \par\vspace{1mm}
}

\begin{document}

\begin{center}
  \LARGE{\textbf{Moteurs électriques --- Corrigé}}
  \vspace{3mm}
\end{center}

\hrule
\vspace{4mm}

%%%%%%%%%%%%%%%%%%%%%%%%%%%%%%%%%%%%%%%%%%
\noindent\textbf{Partie A --- Identification des moteurs} \hfill \textbf{/4 points}

\vspace{3mm}

\Q{Identifier la catégorie de chaque moteur.}

\reponse{
\textbf{Moteur M1 (Leroy-Somer LS80L) :} la plaque indique \textbf{24\,V DC} → alimentation en tension
\emph{continue} → il s'agit d'un \textbf{moteur à courant continu (MCC)}.

\medskip
\textbf{Moteur M2 (Siemens 1LA7) :} la plaque indique \textbf{230/400\,V AC} et une fréquence
\textbf{f = 50\,Hz} → alimentation en tension \emph{alternative} → il s'agit d'un \textbf{moteur asynchrone triphasé}.
}

\Q{Calculer le nombre de paires de pôles, le glissement $g$ du moteur M2.}

\reponse{
Vitesse de synchronisme $n_s = 3000$\,tr/min pour $f = 50$\,Hz :
\[n_s = \frac{60\,f}{p} \implies p = \frac{60 \times 50}{3000} = \mathbf{1} \text{ paire de pôles}\]

Glissement :
\[g = \frac{n_s - n}{n_s} = \frac{3000 - 2840}{3000} = \frac{160}{3000}\]
\[\boxed{g \approx 0{,}053 = 5{,}3~\%}\]
}

\vfill
\newpage

%%%%%%%%%%%%%%%%%%%%%%%%%%%%%%%%%%%%%%%%%%
\noindent\textbf{Partie B --- Bilan de puissance et rendement} \hfill \textbf{/6 points}

\vspace{3mm}

\Q{Calculer la puissance électrique absorbée $P_{\text{élec}}$ par le moteur M1.}

\reponse{
\[P_{\text{élec}} = U \times I = 24 \times 3{,}8\]
\[\boxed{P_{\text{élec}} = 91{,}2~\text{W}}\]
}

\Q{Calculer la vitesse angulaire $\omega$ puis la puissance mécanique utile $P_{\text{méc}}$.}

\reponse{
\[\omega = \frac{2\pi\,n}{60} = \frac{2\pi \times 1420}{60} \approx 148{,}7~\text{rad/s}\]
\[P_{\text{méc}} = C \times \omega = 0{,}55 \times 148{,}7\]
\[\boxed{P_{\text{méc}} \approx 81{,}8~\text{W}}\]
}

\Q{Déduire les pertes $P_{\text{pertes}}$ et le rendement $\eta$.
Comparer à la plaque signalétique et commenter.}

\reponse{
Pertes :
\[P_{\text{pertes}} = P_{\text{élec}} - P_{\text{méc}} = 91{,}2 - 81{,}8 \implies \boxed{P_{\text{pertes}} \approx 9{,}4~\text{W}}\]

Rendement :
\[\eta = \frac{P_{\text{méc}}}{P_{\text{élec}}} = \frac{81{,}8}{91{,}2} \implies \boxed{\eta \approx 0{,}897 = 89{,}7~\%}\]

\textbf{Comparaison avec la plaque :} la plaque indique $P = 90$\,W pour $I_{\rm nom} = 4{,}2$\,A, soit
$P_{\rm élec,nom} = 24 \times 4{,}2 = 100{,}8$\,W, d'où $\eta_{\rm plaque} = 90 / 100{,}8 \approx 89{,}3~\%$.\\[2mm]
\textbf{Commentaire :} les deux valeurs sont très proches ($\approx 89~\%$), ce qui est cohérent. Le moteur
fonctionne légèrement en dessous du courant nominal (3,8\,A $<$ 4,2\,A) → résultat attendu.
}

\vfill
\newpage

%%%%%%%%%%%%%%%%%%%%%%%%%%%%%%%%%%%%%%%%%%
\noindent\textbf{Partie C --- Influence des paramètres électriques sur la vitesse} \hfill \textbf{/10 points}

\vspace{3mm}

\Q{Tracer le graphique $n = f(U)$ et décrire l'évolution.}

\reponse{
\begin{center}
\begin{tikzpicture}
\begin{axis}[
  width=12cm, height=6cm,
  xlabel={Tension $U$ (V)},
  ylabel={Vitesse $n$ (tr/min)},
  xmin=0, xmax=26,
  ymin=0, ymax=1600,
  grid=both,
  grid style={line width=0.3pt, draw=gray!30},
  major grid style={line width=0.6pt, draw=gray!55},
  xtick={0,2,...,26},
  ytick={0,200,...,1600},
  tick label style={font=\small},
  label style={font=\small\bfseries},
  axis line style={thick},
]
  % Points de mesure
  \addplot[only marks, mark=*, mark size=2.5pt, color=blue!70!black]
    coordinates {(4,220)(8,450)(12,680)(16,920)(20,1170)(24,1430)};
  % Droite de tendance (régression linéaire : n ≈ 59.7U - 19)
  \addplot[red!70!black, very thick, domain=0:25] {59.7*x - 19};
\end{axis}
\end{tikzpicture}
\end{center}

\textbf{Description :} la vitesse $n$ est \textbf{proportionnelle} à la tension $U$ (les points sont
alignés sur une droite passant par l'origine). Lorsque $U$ double, $n$ double : c'est une relation
\textbf{linéaire}.
}

\Q{Conséquence pratique pour le contrôle de la vitesse de la perceuse.}

\reponse{
En faisant varier la \textbf{tension d'alimentation} du moteur M1 (par exemple à l'aide d'un variateur
de tension ou d'un hacheur), on contrôle directement la \textbf{vitesse de rotation} de la perceuse.
Cela permet d'adapter la vitesse de coupe selon la dureté du matériau percé.
}

\Q{Tracer le graphique $n = f(f_{\text{réseau}})$ et décrire l'évolution.}

\reponse{
\begin{center}
\begin{tikzpicture}
\begin{axis}[
  width=12cm, height=6cm,
  xlabel={Fréquence du réseau $f$ (Hz)},
  ylabel={Vitesse $n$ (tr/min)},
  xmin=0, xmax=55,
  ymin=0, ymax=3200,
  grid=both,
  grid style={line width=0.3pt, draw=gray!30},
  major grid style={line width=0.6pt, draw=gray!55},
  xtick={0,5,...,55},
  ytick={0,400,...,3200},
  tick label style={font=\small},
  label style={font=\small\bfseries},
  axis line style={thick},
]
  % Points de mesure
  \addplot[only marks, mark=*, mark size=2.5pt, color=blue!70!black]
    coordinates {(10,560)(20,1140)(30,1710)(40,2280)(50,2840)};
  % Droite de tendance (n ≈ 56.8f)
  \addplot[red!70!black, very thick, domain=0:52] {56.8*x};
\end{axis}
\end{tikzpicture}
\end{center}

\textbf{Description :} la vitesse $n$ est \textbf{proportionnelle} à la fréquence $f$ (droite passant
par l'origine). C'est cohérent avec la relation $n_s = \dfrac{60\,f}{p}$ : doubler $f$ double $n_s$
et donc $n$.
}

\Q{Expliquer le principe du variateur de fréquence et donner un exemple d'application.}

\reponse{
\textbf{Principe :} un variateur de fréquence convertit la tension du réseau (50\,Hz fixe) en une
tension alternative de \textbf{fréquence réglable}. Puisque la vitesse du moteur asynchrone est
proportionnelle à la fréquence d'alimentation ($n \propto f$), on règle ainsi la vitesse de rotation
en agissant uniquement sur la fréquence.

\medskip
\textbf{Exemple d'application industrielle :} convoyeur de chaîne de production, pompe centrifuge,
ventilateur industriel, compresseur à débit variable.
}

\vfill

\begin{center}
  \large{\textbf{--- Barème total : /20 points ---}}
\end{center}

\end{document}
